% !TEX root = ../main.tex
\section{Threats to Validity}

We organize validity threats following the taxonomy of Cook and Campbell~\cite{cook1979quasi}, examining internal, external, construct, and statistical conclusion validity.

\subsection{Internal Validity}

The controlled experimental design with automated task execution and fault injection provides strong internal validity. However, several threats remain:

\begin{itemize}
    \item \textbf{Confounding between agent and benchmark:} The full study agents (Claude 3.5 Sonnet, GPT-4o, Gemini 1.5 Pro, Llama 3.3 70B, Qwen 2.5 72B, DeepSeek Chat) were evaluated across all four benchmarks, while Ollama-based agents only ran on two benchmarks (MockBenchmark and GAIA). Benchmark effects cannot be fully disentangled from agent capability differences. Our ablation study confirms that benchmark-specific variance (Task 6) partially confounds cross-agent comparisons.

    \item \textbf{Fault injection ordering:} Faults were injected at fixed rates across runs rather than randomized per trial. Order effects within runs could influence recovery behavior if models adapt to fault patterns over time. The consistent binary recovery pattern (100\% vs. 0\%) suggests this effect is minimal, but counterbalanced injection sequences would strengthen causal attribution.

    \item \textbf{Seed-based variability across metrics:} Sensitivity analysis (Task 3) reveals that 60\% of agents show zero seed variance, but ollama models exhibit rank changes of up to 2 positions across the 1,262 unique seeds examined. Uncontrolled seed variation across experiments introduces uncontrolled variance in reliability estimates for these models.
\end{itemize}

\subsection{External Validity}

\begin{itemize}
    \item \textbf{Task generalizability:} The evaluation uses simulated (mock) tasks for the majority of evaluations (27,160 of 31,981). While this enables controlled fault injection and perturbation testing, simulated tasks may not capture the complexity and variability of real-world agentic workloads. The GAIA and SWEBench Lite evaluations provide some external grounding, but these benchmarks themselves have known limitations in ecological validity~\cite{mialon2023gaia, jimenez2024swebench}.

    \item \textbf{Fault type coverage:} Six fault types were injected, but production systems encounter additional failure modes including rate limiting, authentication token expiration, model degradation under sustained load, and gradual drift in API response quality. Our finding that network interruptions and tool failures are catastrophic (0\% recovery) may not generalize to fault implementations with different retry or fallback strategies.

    \item \textbf{Model selection and recency:} The 21 evaluated agents represent a convenience sample of publicly available models as of mid-2025. New model releases, API updates, and quantization changes may alter reliability characteristics. The pronounced ceiling effect limits generalizability to future, more capable models.
\end{itemize}

\subsection{Construct Validity}

\begin{itemize}
    \item \textbf{Metric saturation:} The composite reliability metric is defined as the unweighted mean of success rate, repeated-run consistency, fault tolerance, and perturbation robustness. For 16 of 21 agents, this score saturates at 1.0. This ceiling effect (documented in Tasks 1, 3, and 4) indicates either that the metric lacks sufficient resolution or that current LLMs genuinely achieve maximal reliability under our evaluation conditions. Our ablation study confirms that this saturation is not an artifact of any single component metric.

    \item \textbf{Operationalization of reliability:} Our definition of reliability focuses on four observable dimensions. Other conceptions of reliability -- such as calibration quality~\cite{guo2017calibration}, output factuality consistency~\cite{min2023factual}, or decision boundary stability~\cite{jiang2021can} -- are not captured. The strong agreement across our four dimensions (inter-metric correlation $\geq 0.85$ for the full-metric group) suggests internal consistency but does not guarantee coverage of the broader construct.
\end{itemize}

\subsection{Statistical Conclusion Validity}

\begin{itemize}
    \item \textbf{Limited sample size for differentiation:} With 3--9 repeated runs per agent-benchmark pair for API-based models, the statistical power to detect small effect sizes ($d < 0.5$) between top-tier models is limited. Bootstrap confidence intervals (Task 4) show that pairwise comparisons among the six top models yield no statistically significant differences, though effect sizes for top-vs-lightweight model comparisons are large (Cohen's $d > 2.0$).

    \item \textbf{Missing data and null metrics:} Ollama-based agents lack fault tolerance and perturbation robustness metrics, reducing them to a 2-metric composite versus the 4-metric composite for API-based agents. This structural missingness biases cross-group comparisons. Our weighting sensitivity analysis (Task 2) shows that Ollama rankings are more sensitive to weighting scheme choice (rank change $\sigma = 1.2$ versus $\sigma = 0.0$ for full-metric agents).

    \item \textbf{Correlation of ranking metrics:} The cross-metric correlation matrix (Table~\ref{tab:metric_correlation}) reveals that success rate and repeated-run consistency are perfectly correlated ($\rho = 1.0$) in the Ollama group, indicating redundancy. The ablation study confirms that removing either metric individually has identical effects on ranking for these agents.
\end{itemize}
